\subsection{Subgroupoids}

DEFINITION 5. --- Let C be a category. A subcategory of C is a subquiver D of C satisfying the following conditions:

\begin{enumerate}
\def\labelenumi{(\roman{enumi})}
\item
  For every vertex a of D, \(e_a\) is an arrow of D;
\item
  For every composable pair \((f,g)\) of arrows of C belonging to D, the product fg is an arrow of D;
\end{enumerate}

Let C be a category and let D be a subcategory of C; endowed with the composition law induced from that of C by restriction to subsets, D is a category.

Every full subquiver of C is a subcategory of C. The intersection of a family of subcategories of C is a subcategory of C.

DEFINITION 6. --- Let G be a groupoid. A subgroupoid of G is a subcategory H of G such that the inverse of every arrow of G belonging to H is an arrow of H.

Let G be a groupoid. Every subgroupoid of G is a groupoid. Every full subquiver of G is a subgroupoid of G. The intersection of a family of subgroupoids of G is a subgroupoid of G.

Let H be a subquiver of G. The intersection of the subgroupoids of G of which H is a subquiver is called the subgroupoid of G generated by H. The set of its vertices is equal to that of H and its orbits are the connected components of H.

Examples. --- 1) Let X be a set; denote by X × X the groupoid of pairs of X (II, p. 163, example 2). Let R be an equivalence relation on X. The subquiver of the groupoid X × X whose set of vertices is X and whose set of arrows is the graph of the equivalence relation R is a subgroupoid of X × X. Its orbits are the equivalence classes of the relation R.

Conversely, every subgroupoid of X × X whose set of vertices is X is of this form.

Let X and S be sets and let \(p: X \rightarrow S\) be a map. The map p defines an equivalence relation on X (E, II, p. 41). We denote by \(X \times_{S} X\) the subgroupoid of \(X \times X\) defined by this equivalence relation and call it the groupoid of pairs of \((X, p)\) .

\begin{enumerate}
\def\labelenumi{\arabic{enumi})}
\setcounter{enumi}{1}
\item
  Let G and G' be groupoids, and let \(\varphi\) and \(\psi\) be morphisms of groupoids from G to G'. The equalizer of \(\varphi\) and \(\psi\) (cf. II, p. 153) is a subgroupoid of G.
\item
  Let X be a set and \(\Gamma\) a group. Denote by \(X \times \Gamma \times X\) the quiver whose set of vertices is X, whose set of arrows is \(X \times \Gamma \times X\) , the source and target maps being respectively \(\mathrm{pr}_1\) and \(\mathrm{pr}_3\) . Endowed with the composition law \((x, \gamma, x') \cdot (x', \gamma', x'') = (x, \gamma \gamma', x'')\) , this is a groupoid. It is moreover the groupoid deduced by transport of structure from the groupoid \((\mathrm{X}\times\mathrm{X})\times\Gamma\) , the product of the groupoid of pairs \(X\times X\) and the groupoid associated with the group \(\Gamma\) (II, p. 163, example 1), by means of the bijection \((x,x',\gamma)\mapsto(x,\gamma,x')\) from \(X\times X\times\Gamma\) to \(X\times\Gamma\times X\) .
\end{enumerate}

Let \(m\colon X \times \Gamma \to X\) be an operation (on the right) of the group \(\Gamma\) on \(X\) . The graph \(F\) of the map \(m\) , that is, the set of triples \((x, \gamma, x') \in X \times \Gamma \times X\) such that \(x' = x\gamma\) , is the set of arrows of a unique subgroupoid \(G\) of \(X \times \Gamma \times X\) . The orbits of this groupoid are the orbits of \(\Gamma\) on \(X\) ; if \(x \in X\) , the isotropy group of \(G\) at \(x\) is the set of \((x, \gamma, x)\) , where \(\gamma\) ranges over the stabilizer of \(x\) in \(\Gamma\) .

Conversely, any subgroupoid G of the groupoid \(X \times \Gamma \times X\) such that the map \((\mathrm{pr}_{1}, \mathrm{pr}_{2})\) from \(\mathrm{Fl}(\mathrm{G})\) to \(X \times \Gamma\) is a bijection is of this form.

\begin{enumerate}
\def\labelenumi{\arabic{enumi})}
\setcounter{enumi}{3}
\tightlist
\item
  Let G be a groupoid, let X be a set, and let \(\varphi\colon\mathrm{X}\to\mathrm{Som}(\mathrm{G})\) be a map. We define a quiver \(\mathrm{G}'\) as follows. The set \(\mathrm{Som}(\mathrm{G}')\) is the set X; for every pair \((x,y)\) of elements of X, the set \(\mathrm{Fl}_{x,y}(\mathrm{G}')\) is the set of triples \((x,f,y)\in\mathrm{X}\times\mathrm{Fl}(\mathrm{G})\times\mathrm{X}\) , where \(f\) is an element of \(\mathrm{Fl}_{\varphi(x),\varphi(y)}(\mathrm{G})\) . Let \(x,y,z\) be elements of X, \(f\in\mathrm{Fl}_{\varphi(x),\varphi(y)}(\mathrm{G})\) and \(g\in\mathrm{Fl}_{\varphi(y),\varphi(z)}(\mathrm{G})\) be arrows of G, and set \((x,f,y)\cdot(y,g,z)=(x,fg,z)\) . This defines a composition law in the quiver \(\mathrm{G}'\), making it a groupoid. It is called the inverse image groupoid of the groupoid G by the map \(\varphi\) and we denote it \(\varphi^{*}(\mathrm{G})\) .
\end{enumerate}

The pair \((\varphi,\psi)\) , where \(\psi\colon\mathrm{Fl}(\varphi^{*}(\mathrm{G}))\to\mathrm{Fl}(\mathrm{G})\) is the map defined by \((x,f,y)\mapsto f\) , is a morphism of groupoids from \(\varphi^{*}(\mathrm{G})\) to G, called the canonical morphism.

\begin{enumerate}
\def\labelenumi{\arabic{enumi})}
\setcounter{enumi}{4}
\tightlist
\item
  Let \(\varphi\colon\mathrm{G}\to\mathrm{G}^{\prime}\) be a morphism of groupoids. Let H be the subquiver of G with vertex set \(\mathrm{Som}(\mathrm{G})\) and whose set of arrows consists of the \(f\in\mathrm{Fl}(\mathrm{G})\) such that \(\varphi(f)\) is an identity element of \(\mathrm{G}^{\prime}\) . For all \(a\in\mathrm{Som}(\mathrm{G})\) , \(e_{a}\in\mathrm{H}_{a,a}\) . For all \(f\in\mathrm{Fl}(\mathrm{H})\) , \(\varphi(f^{-1})=\varphi(f)^{-1}\) ; therefore \(f^{-1}\in\mathrm{Fl}(\mathrm{H})\) . Moreover, if \(f\) and \(g\) are composable arrows of H, then \(\varphi(fg)=\varphi(f)\varphi(g)=e_{t(\varphi(f))}e_{o(\varphi(g))}=e_{t(\varphi(f))}\) , therefore \(fg\) is an arrow of H. This proves that H is a subgroupoid of G. It is called the kernel of \(\varphi\) and we denote it \(\mathrm{Ker}(\varphi)\) .
\end{enumerate}

\pandocbounded{\includegraphics[keepaspectratio]{images/part002_cf856079.jpg}}

\begin{enumerate}
\def\labelenumi{\arabic{enumi})}
\setcounter{enumi}{5}
\tightlist
\item
  Let \(\varphi\colon\mathrm{G}\to\mathrm{G}^{\prime}\) be a morphism of groupoids. The quiver \(\varphi(\mathrm{G})\) whose set of vertices is \(\varphi(\mathrm{Som}(\mathrm{G}))\) and whose set of arrows is \(\varphi(\mathrm{Fl}(\mathrm{G}))\) is in general not a subgroupoid of \(\mathrm{G}^{\prime}\) .
\end{enumerate}

This is nevertheless the case if the map Som( \(\varphi\) ) is injective (II, p. 225, exerc. 5).

